% ============================================================
%  RAPPORT DE FIN D'ETUDES
%  Module : Authentification et Gestion des Roles
%  Etudiant  : Ghoulam Mohamed Said
%  Equipe    : Reddas Belkacem, Chelig Mohamed Amine
%  Encadreur : Dr. Abdelkader Bougassa
%  Universite Ibn Khaldoun - Tiaret - 2024/2025
%  Compatible Overleaf - pdfLaTeX
% ============================================================

\documentclass[12pt,a4paper]{report}

\usepackage[utf8]{inputenc}
\usepackage[T1]{fontenc}
\usepackage[french]{babel}
\usepackage{graphicx}
\usepackage{geometry}
\usepackage{booktabs}
\usepackage{xcolor}
\usepackage{hyperref}
\usepackage{float}

\geometry{top=2.5cm,bottom=2.5cm,left=3cm,right=2.5cm}

\definecolor{myblue}{RGB}{30,80,160}
\definecolor{mygold}{RGB}{180,140,30}

\hypersetup{
  colorlinks=true,
  linkcolor=myblue,
  urlcolor=mygold,
  citecolor=myblue
}

\title{Systeme de Gestion Universitaire Integre \\
       Module : Authentification et Gestion des Roles}
\author{Ghoulam Mohamed Said}
\date{2024 -- 2025}

% ============================================================
\begin{document}
% ============================================================

% ────────────────────────────────────────────────────────────
%  PAGE DE GARDE
% ────────────────────────────────────────────────────────────
\begin{titlepage}
\begin{center}

\textsc{\large Republique Algerienne Democratique et Populaire}\\[0.2cm]
\textsc{Ministere de l'Enseignement Superieur et de la Recherche Scientifique}\\[0.5cm]

\rule{\linewidth}{0.5mm}\\[0.3cm]
\textsc{\Large Universite Ibn Khaldoun -- Tiaret}\\[0.2cm]
\textsc{Faculte des Mathematiques et de l'Informatique}\\[0.2cm]
\textsc{Departement d'Informatique}\\
\rule{\linewidth}{0.5mm}\\[1.5cm]

{\large Memoire de Fin d'Etudes}\\[0.3cm]
{\large En vue de l'obtention du diplome de Master en Informatique}\\[0.2cm]
{\large Specialite : Genie Logiciel (ISIL)}\\[2cm]

\textcolor{myblue}{\rule{0.6\linewidth}{1pt}}\\[0.5cm]
{\Huge \bfseries Systeme de Gestion}\\[0.3cm]
{\Huge \bfseries Universitaire Integre}\\[0.5cm]
{\Large \textcolor{mygold}{Module : Authentification et Gestion des Roles}}\\[0.5cm]
\textcolor{myblue}{\rule{0.6\linewidth}{1pt}}\\[2cm]

\begin{minipage}{0.45\textwidth}
\begin{flushleft}
\textbf{Realise par :}\\[0.2cm]
Ghoulam Mohamed Said\\
Reddas Belkacem\\
Chelig Mohamed Amine
\end{flushleft}
\end{minipage}
\hfill
\begin{minipage}{0.45\textwidth}
\begin{flushright}
\textbf{Encadre par :}\\[0.2cm]
Dr. Abdelkader Bougassa\\
\textit{Maitre de Conferences}\\
Universite Ibn Khaldoun
\end{flushright}
\end{minipage}\\[2cm]

{\large \textbf{Annee Universitaire : 2024 -- 2025}}

\end{center}
\end{titlepage}

% ────────────────────────────────────────────────────────────
%  DEDICACE
% ────────────────────────────────────────────────────────────
\chapter*{Dedicace}
\addcontentsline{toc}{chapter}{Dedicace}

\vspace*{2cm}
\begin{center}
\begin{minipage}{0.85\textwidth}
\begin{center}
\textit{Je dedie ce modeste travail :}\\[0.5cm]

\textit{A mes chers parents, source inepuisable d'amour, de patience}\\
\textit{et de sacrifices, qui m'ont toujours soutenu durant mon parcours.}\\[0.5cm]

\textit{A mes freres et sceurs pour leur soutien inconditionnel.}\\[0.5cm]

\textit{A tous mes enseignants qui ont eclaire mon chemin.}\\[0.5cm]

\textit{A mes camarades de promotion et a tous ceux}\\
\textit{qui croient en la valeur de la connaissance.}\\[1cm]

\textbf{-- Ghoulam Mohamed Said --}
\end{center}
\end{minipage}
\end{center}

% ────────────────────────────────────────────────────────────
%  REMERCIEMENTS
% ────────────────────────────────────────────────────────────
\chapter*{Remerciements}
\addcontentsline{toc}{chapter}{Remerciements}

Au terme de ce travail, nous tenons a exprimer notre profonde gratitude
a notre encadreur, \textbf{Dr. Abdelkader Bougassa}, pour ses precieux
conseils, sa disponibilite constante et la rigueur scientifique dont il
a fait preuve tout au long de l'encadrement de ce projet.

Nos sinceres remerciements vont egalement aux membres du jury qui ont
accepte d'evaluer ce travail, ainsi qu'a l'ensemble du corps enseignant
du departement d'informatique pour la qualite de la formation dispensee
durant notre cursus universitaire.

Nous remercions enfin nos familles respectives pour leur soutien moral
et leur patience, et nos camarades pour les echanges enrichissants qui
ont contribue a l'aboutissement de ce projet.

% ────────────────────────────────────────────────────────────
%  RESUME
% ────────────────────────────────────────────────────────────
\chapter*{Resume}
\addcontentsline{toc}{chapter}{Resume}

Ce memoire presente la conception et le developpement d'un \textbf{Systeme
de Gestion Universitaire Integre} (SGUI), centre sur le module
\textbf{d'Authentification et de Gestion des Roles}. Le systeme vise a
numeriser les processus administratifs d'un etablissement d'enseignement
superieur algerien en offrant une plateforme unifiee accessible a trois
profils d'utilisateurs distincts : Administrateur, Enseignant et Etudiant.

Le module developpe implemente un mecanisme d'authentification securise
base sur les \textbf{JSON Web Tokens (JWT)}, un controle d'acces a base
de roles (\textbf{RBAC}), ainsi qu'une gestion complete du cycle de vie
des sessions. Le backend repose sur \textbf{Node.js} avec \textbf{TypeScript},
exposant une API RESTful construite avec \textbf{Express.js}, tandis que
la persistance des donnees est assuree par \textbf{PostgreSQL} via l'ORM
\textbf{Prisma}. L'interface utilisateur est developpee en
\textbf{React.js} avec le framework CSS \textbf{Tailwind CSS}.

Les resultats obtenus demontrent la robustesse du systeme face aux
principales vulnerabilites de securite (injection SQL, XSS, CSRF,
brute-force) et confirment sa capacite a gerer efficacement les flux
d'acces multi-roles dans un contexte universitaire reel.

\vspace{0.5cm}
\textbf{Mots-cles :} Authentification, JWT, RBAC, API REST, Node.js,
React.js, PostgreSQL, Prisma, Securite web, Gestion universitaire.

% ────────────────────────────────────────────────────────────
%  ABSTRACT
% ────────────────────────────────────────────────────────────
\chapter*{Abstract}
\addcontentsline{toc}{chapter}{Abstract}

This thesis presents the design and development of an \textbf{Integrated
University Management System} (IUMS), focusing on the \textbf{Authentication
and Role Management} module. The system aims to digitize the administrative
processes of an Algerian higher education institution by providing a
unified platform accessible to three distinct user profiles: Administrator,
Teacher, and Student.

The developed module implements a secure authentication mechanism based on
\textbf{JSON Web Tokens (JWT)}, \textbf{Role-Based Access Control (RBAC)},
and full session lifecycle management. The backend is built with
\textbf{Node.js} and \textbf{TypeScript}, exposing a RESTful API via
\textbf{Express.js}, with data persistence handled by \textbf{PostgreSQL}
through the \textbf{Prisma} ORM. The frontend is developed in
\textbf{React.js} with \textbf{Tailwind CSS}.

Experimental results demonstrate the system's robustness against the main
web security vulnerabilities (SQL injection, XSS, CSRF, brute-force) and
confirm its ability to efficiently manage multi-role access flows in a
real university context.

\vspace{0.5cm}
\textbf{Keywords:} Authentication, JWT, RBAC, REST API, Node.js, React.js,
PostgreSQL, Prisma, Web security, University management.

% ────────────────────────────────────────────────────────────
%  TABLE DES MATIERES
% ────────────────────────────────────────────────────────────
\tableofcontents
\listoffigures
\listoftables

% ════════════════════════════════════════════════════════════
%  INTRODUCTION GENERALE
% ════════════════════════════════════════════════════════════
\chapter*{Introduction Generale}
\addcontentsline{toc}{chapter}{Introduction Generale}

\section*{Contexte et Motivation}

La transformation numerique des etablissements d'enseignement superieur
constitue aujourd'hui un enjeu strategique majeur. En Algerie, de nombreuses
universites gerent encore leurs processus administratifs de maniere
fragmentee : formulaires papier, tableurs disperses, communications
informelles par messagerie. Cette situation genere des inefficacites, des
risques d'erreur et une experience utilisateur degradee pour les etudiants
comme pour le personnel academique.

Le present projet s'inscrit dans ce contexte en proposant un \textbf{Systeme
de Gestion Universitaire Integre (SGUI)} qui centralise et numerise les
principaux processus : gestion des projets de fin d'etudes (PFE),
reclamations et justifications d'absence, annonces institutionnelles,
remise de copies, et gestion du jury.

\section*{Problematique}

La question centrale de ce travail est la suivante : comment concevoir un
systeme d'authentification et de controle d'acces robuste, evolutif et
convivial, capable de gerer des profils utilisateurs aux privileges
heterogenes dans un environnement universitaire reel ?

Cette problematique souleve plusieurs defis techniques :
\begin{itemize}
  \item Garantir la confidentialite et l'integrite des donnees sans
        introduire de frictions excessives pour l'utilisateur.
  \item Implementer un modele de controle d'acces suffisamment fin pour
        distinguer administrateurs, enseignants et etudiants.
  \item Assurer la scalabilite du systeme face a une population
        universitaire potentiellement importante.
  \item Proteger le systeme contre les attaques courantes du web
        (injection, XSS, CSRF, force brute).
\end{itemize}

\section*{Objectifs}

Les objectifs de ce memoire sont les suivants :
\begin{enumerate}
  \item Analyser l'existant et identifier les besoins fonctionnels et non
        fonctionnels du SGUI.
  \item Concevoir l'architecture logicielle du systeme avec un accent
        particulier sur le module d'authentification.
  \item Developper et tester un module d'authentification JWT securise
        avec RBAC a trois roles.
  \item Valider le systeme par des tests fonctionnels et de securite.
\end{enumerate}

\section*{Structure du Memoire}

Ce document est organise en trois chapitres :
\begin{itemize}
  \item \textbf{Chapitre 1} -- Etat de l'art : revue bibliographique des
        technologies et concepts utilises (JWT, RBAC, API REST,
        frameworks modernes).
  \item \textbf{Chapitre 2} -- Analyse et conception : specification des
        besoins, modelisation UML et architecture globale.
  \item \textbf{Chapitre 3} -- Implementation et tests : detail du
        developpement du module d'authentification, suivi d'une campagne
        de tests fonctionnels et de securite.
\end{itemize}

% ════════════════════════════════════════════════════════════
\chapter{Etat de l'Art et Etude du Domaine}
% ════════════════════════════════════════════════════════════

\section{Introduction}

Ce chapitre presente les fondements theoriques et technologiques sur
lesquels repose le module d'authentification et de gestion des roles du
SGUI. Nous passons en revue les concepts de securite applicative, les
paradigmes d'authentification modernes, et les technologies choisies pour
leur pertinence et leur maturite dans l'ecosysteme du developpement web.

\section{Concepts Fondamentaux de Securite Web}

\subsection{Authentification vs. Autorisation}

L'\textbf{authentification} est le processus permettant de verifier
l'identite d'un utilisateur (prouver qui il est). L'\textbf{autorisation}
est le processus determinant ce qu'un utilisateur authentifie est
autorise a faire (definir ce qu'il peut faire).

Ces deux concepts sont fondamentalement distincts mais complementaires.
L'authentification precede toujours l'autorisation dans le flux de
traitement d'une requete.

\subsection{Les Principaux Mecanismes d'Authentification Web}

\subsubsection{Sessions cote serveur}

Le mecanisme traditionnel repose sur des sessions stockees cote serveur.
A la connexion, le serveur cree une session identifiee par un identifiant
unique transmis au client via un cookie. Ce modele presente l'inconvenient
d'introduire un etat cote serveur, limitant la scalabilite horizontale.

\subsubsection{JSON Web Tokens (JWT)}

Le standard \textbf{RFC 7519} definit le JWT comme un moyen compact et
auto-contenu de transmettre des informations entre parties sous forme d'un
objet JSON signe. Un JWT est compose de trois parties encodees en
Base64Url et separees par des points :

\begin{itemize}
  \item \textbf{Header} : algorithme de signature (HS256, RS256) et type
        du token.
  \item \textbf{Payload} : claims (affirmations) sur l'utilisateur et
        metadonnees (sub, role, exp, iat).
  \item \textbf{Signature} : garantit l'integrite du token.
\end{itemize}

L'avantage majeur du JWT est son caractere \textbf{stateless} : le
serveur n'a pas besoin de stocker l'etat de session, ce qui facilite le
passage a l'echelle.

\subsection{Controle d'Acces a Base de Roles (RBAC)}

Le RBAC est un modele de controle d'acces dans lequel les permissions
sont associees a des \textbf{roles}, et les utilisateurs sont assignes
a ces roles. Un utilisateur herite ainsi des permissions de tous ses
roles assignes.

\begin{table}[H]
\centering
\caption{Entites du modele RBAC}
\begin{tabular}{lll}
\toprule
\textbf{Entite} & \textbf{Description} & \textbf{Exemple} \\
\midrule
Utilisateur & Sujet humain ou systeme & Etudiant, Enseignant \\
Role        & Ensemble de permissions & admin, teacher, student \\
Permission  & Droit sur une ressource & WRITE:/pfe/sujets \\
\bottomrule
\end{tabular}
\end{table}

\section{Technologies et Frameworks}

\subsection{Backend -- Node.js et TypeScript}

\textbf{Node.js} est un environnement d'execution JavaScript cote serveur
base sur le moteur V8 de Chrome. Son modele d'entree/sortie non-bloquant
le rend particulierement adapte aux API RESTful a haute concurrence.

\textbf{TypeScript} est un sur-ensemble type de JavaScript qui ameliore
la maintenabilite du code a travers la verification statique des types,
reduisant les erreurs a l'execution.

\textbf{Express.js} est un framework minimaliste pour Node.js offrant un
systeme de routage et de middlewares flexible.

\subsection{ORM -- Prisma}

\textbf{Prisma} est un ORM de nouvelle generation pour Node.js et
TypeScript. Il offre :
\begin{itemize}
  \item Un schema declaratif comme source de verite.
  \item Un client type genere automatiquement.
  \item Des migrations versionnees.
  \item Une protection native contre les injections SQL.
\end{itemize}

\subsection{Base de Donnees -- PostgreSQL}

\textbf{PostgreSQL} est un SGBD relationnel open-source repute pour sa
conformite aux standards SQL, sa robustesse et ses capacites avancees
(types JSON, contraintes, triggers).

\subsection{Frontend -- React.js et Tailwind CSS}

\textbf{React.js} est une bibliotheque JavaScript developpee par Meta
pour la construction d'interfaces utilisateur basees sur des composants
reutilisables.

\textbf{Tailwind CSS} est un framework CSS utilitaire permettant de
concevoir des interfaces modernes directement dans le JSX.

\subsection{Comparaison des Solutions d'Authentification}

\begin{table}[H]
\centering
\caption{Comparaison des mecanismes d'authentification}
\begin{tabular}{lcccc}
\toprule
\textbf{Critere}    & \textbf{Session} & \textbf{JWT} & \textbf{OAuth2} & \textbf{API Key} \\
\midrule
Stateless           & Non       & Oui         & Partiel & Oui          \\
Scalabilite         & Faible    & Elevee      & Elevee  & Elevee       \\
Revocation          & Facile    & Complexe    & Standard & Facile      \\
Complexite          & Faible    & Moyenne     & Elevee  & Tres faible  \\
Adapte au SGUI      & Possible  & \textbf{Oui}& Excessif& Non          \\
\bottomrule
\end{tabular}
\end{table}

\section{Conclusion}

Ce chapitre a pose les bases conceptuelles et technologiques du projet.
Nous avons justifie le choix du JWT pour l'authentification et du RBAC
pour la gestion des acces, et presente la stack technique retenue. Le
chapitre suivant detaille l'analyse des besoins et la conception du
systeme.

% ════════════════════════════════════════════════════════════
\chapter{Analyse et Conception}
% ════════════════════════════════════════════════════════════

\section{Introduction}

Ce chapitre presente la phase d'analyse et de conception du SGUI. Nous
definissons d'abord les acteurs et leurs besoins, puis nous modelisons
le systeme a l'aide de diagrammes UML couvrant les cas d'utilisation,
les sequences d'interaction et la structure statique des donnees.

\section{Analyse des Besoins}

\subsection{Identification des Acteurs}

\begin{table}[H]
\centering
\caption{Acteurs du systeme et leurs responsabilites}
\begin{tabular}{lp{4cm}p{6cm}}
\toprule
\textbf{Acteur} & \textbf{Role} & \textbf{Responsabilites} \\
\midrule
Administrateur & Gestionnaire de la plateforme &
  Creer des comptes, gerer les PFE, valider les reclamations,
  publier des annonces \\
\addlinespace
Enseignant & Personnel academique &
  Proposer des sujets PFE, s'auto-inscrire au jury, traiter les
  justifications \\
\addlinespace
Etudiant & Apprenant &
  Consulter les annonces, soumettre des reclamations, deposer des
  travaux \\
\bottomrule
\end{tabular}
\end{table}

\subsection{Besoins Fonctionnels}

\subsubsection{Module Authentification et Roles}

\begin{enumerate}
  \item L'utilisateur doit pouvoir s'authentifier via un identifiant et
        un mot de passe.
  \item Le systeme doit emettre un token JWT signe a la connexion reussie.
  \item Le systeme doit refuser l'acces aux ressources protegees si le
        token est absent, expire ou invalide.
  \item Chaque route API doit verifier que le role de l'utilisateur
        dispose des permissions requises.
  \item Un administrateur doit pouvoir creer, modifier et desactiver des
        comptes utilisateurs.
  \item Le systeme doit permettre la deconnexion explicite.
  \item Un mecanisme de rafraichissement de token doit etre disponible.
\end{enumerate}

\subsection{Besoins Non Fonctionnels}

\begin{table}[H]
\centering
\caption{Besoins non fonctionnels}
\begin{tabular}{lp{9cm}}
\toprule
\textbf{Categorie} & \textbf{Exigence} \\
\midrule
Securite       & Mots de passe haches (bcrypt cout 12), tokens
                 expirables, protection CORS, rate-limiting \\
Performance    & Temps de reponse inferieur a 500 ms pour 95\% des
                 requetes \\
Disponibilite  & Uptime de 99\% hors maintenance planifiee \\
Maintenabilite & Architecture modulaire, TypeScript strict, tests
                 unitaires \\
Portabilite    & Conteneurisable via Docker, independant de l'OS \\
\bottomrule
\end{tabular}
\end{table}

\section{Diagrammes UML}

\subsection{Diagramme de Cas d'Utilisation}

La figure ci-dessous presente les principaux cas d'utilisation du module
Authentification et Gestion des Roles. Trois acteurs interagissent avec
le systeme : Administrateur, Enseignant et Etudiant.

\begin{figure}[H]
\centering
\framebox[0.85\linewidth]{
  \begin{minipage}{0.8\linewidth}
    \vspace{2cm}
    \begin{center}
    \textbf{[ Diagramme de Cas d'Utilisation ]}\\[0.5cm]
    Acteurs : Administrateur, Enseignant, Etudiant\\[0.3cm]
    Cas d'utilisation principaux :\\
    -- S'authentifier\\
    -- Se deconnecter\\
    -- Consulter son profil\\
    -- Gerer les comptes (Admin)\\
    -- Gerer les sujets PFE\\
    -- S'affecter au jury\\
    -- Soumettre une reclamation\\[0.5cm]
    \textit{(A inserer : capture de l'outil de modelisation UML)}
    \end{center}
    \vspace{2cm}
  \end{minipage}
}
\caption{Diagramme de cas d'utilisation -- Module Auth et Roles}
\label{fig:usecase}
\end{figure}

\subsection{Diagramme de Sequence -- Flux d'Authentification}

La figure suivante illustre le flux complet de connexion d'un
utilisateur, de la saisie des identifiants jusqu'a l'emission du JWT.

\begin{figure}[H]
\centering
\framebox[0.85\linewidth]{
  \begin{minipage}{0.8\linewidth}
    \vspace{2cm}
    \begin{center}
    \textbf{[ Diagramme de Sequence -- Authentification JWT ]}\\[0.5cm]
    Lifelines :\\
    Client (React) -- Auth Middleware -- Auth Controller -- PostgreSQL\\[0.5cm]
    Etapes principales :\\
    1. POST /api/auth/login (email, password)\\
    2. findUser(email) sur la base de donnees\\
    3. bcrypt.compare(password, hash)\\
    4. jwt.sign(\{id, role\}, SECRET)\\
    5. Retour 200 \{token, user\}\\
    6. Stockage dans localStorage\\
    7. Authorization: Bearer token\\[0.5cm]
    \textit{(A inserer : capture du diagramme de sequence)}
    \end{center}
    \vspace{2cm}
  \end{minipage}
}
\caption{Diagramme de sequence -- Flux d'authentification JWT}
\label{fig:sequence}
\end{figure}

\subsection{Diagramme de Classes Simplifie}

\begin{figure}[H]
\centering
\framebox[0.85\linewidth]{
  \begin{minipage}{0.8\linewidth}
    \vspace{1.5cm}
    \begin{center}
    \textbf{[ Diagramme de Classes -- Module Auth ]}\\[0.5cm]
    Classes principales :\\[0.3cm]

    \textbf{User}\\
    -- id : String (UUID)\\
    -- email : String (unique)\\
    -- passwordHash : String\\
    -- role : Role (enum)\\
    -- nom, prenom : String\\
    -- isActive : Boolean\\[0.3cm]

    \textbf{Role (enumeration)}\\
    ADMIN, TEACHER, STUDENT\\[0.3cm]

    \textbf{AuthService}\\
    + login(email, pwd)\\
    + verifyToken(token)\\
    + hashPassword(pwd)\\
    + generateToken(user)\\[0.5cm]
    \textit{(A inserer : capture du diagramme de classes)}
    \end{center}
    \vspace{1.5cm}
  \end{minipage}
}
\caption{Diagramme de classes -- Module Authentification}
\label{fig:classes}
\end{figure}

\section{Architecture du Systeme}

Le SGUI suit une architecture Client-Serveur a trois couches :

\begin{table}[H]
\centering
\caption{Architecture en trois couches}
\begin{tabular}{ll}
\toprule
\textbf{Couche} & \textbf{Technologies} \\
\midrule
Presentation        & React.js + Tailwind CSS + Axios \\
Applicative (API)   & Node.js + Express.js + TypeScript \\
Persistance         & PostgreSQL + Prisma ORM \\
\bottomrule
\end{tabular}
\end{table}

\section{Conclusion}

Ce chapitre a presente l'analyse des besoins fonctionnels et non
fonctionnels du SGUI, modelise les interactions cles via des diagrammes
UML et defini l'architecture logicielle retenue. Le chapitre suivant
detaille l'implementation concrete du module d'authentification et de
gestion des roles.

% ════════════════════════════════════════════════════════════
\chapter{Implementation et Tests}
% ════════════════════════════════════════════════════════════

\section{Introduction}

Ce chapitre decrit l'implementation du module Authentification et
Gestion des Roles du SGUI. Nous presentons successivement
l'environnement de developpement, les composants backend (service JWT,
middlewares, controleur), les composants frontend (contexte
d'authentification, routes protegees), puis les resultats des tests
fonctionnels et de securite.

\section{Environnement de Developpement}

\begin{table}[H]
\centering
\caption{Environnement de developpement}
\begin{tabular}{ll}
\toprule
\textbf{Composant} & \textbf{Version} \\
\midrule
Node.js              & 20.x LTS \\
TypeScript           & 5.x \\
Express.js           & 4.18 \\
Prisma               & 5.x \\
PostgreSQL           & 16.x \\
React.js             & 18.x \\
Tailwind CSS         & 3.x \\
bcryptjs             & 2.4.x (cout 12) \\
jsonwebtoken         & 9.x (HS256) \\
\bottomrule
\end{tabular}
\end{table}

\section{Architecture du Module}

Le module suit une architecture en couches. Chaque requete traverse
obligatoirement les middlewares d'authentification et d'autorisation
avant d'atteindre le controleur metier.

Le flux de traitement est le suivant :
\begin{enumerate}
  \item Le client React envoie une requete avec un en-tete
        \texttt{Authorization: Bearer <token>}.
  \item Le middleware \texttt{requireAuth} verifie la signature et
        l'expiration du token.
  \item Le middleware \texttt{requireRole} verifie que le role est
        autorise pour la route.
  \item Le controleur execute la logique metier et interroge la base de
        donnees via Prisma.
\end{enumerate}

\section{Modele de Donnees}

\begin{table}[H]
\centering
\caption{Schema de l'entite User}
\begin{tabular}{lll}
\toprule
\textbf{Champ} & \textbf{Type} & \textbf{Description} \\
\midrule
id           & String (UUID)    & Identifiant unique           \\
email        & String (unique)  & Adresse e-mail               \\
passwordHash & String           & Hash bcrypt du mot de passe  \\
role         & Enum             & admin / teacher / student    \\
nom          & String           & Nom de famille               \\
prenom       & String           & Prenom                       \\
isActive     & Boolean          & Compte actif ou desactive    \\
createdAt    & DateTime         & Date de creation             \\
updatedAt    & DateTime         & Derniere modification        \\
\bottomrule
\end{tabular}
\end{table}

\subsection{Matrice d'Autorisation}

\begin{table}[H]
\centering
\caption{Matrice d'autorisation des routes API}
\begin{tabular}{llccc}
\toprule
\textbf{Methode} & \textbf{Route} & \textbf{Admin} & \textbf{Teacher} & \textbf{Student} \\
\midrule
POST & /api/auth/login            & Oui & Oui & Oui \\
GET  & /api/auth/me               & Oui & Oui & Oui \\
POST & /api/auth/logout           & Oui & Oui & Oui \\
GET  & /api/users                 & Oui & Non & Non \\
POST & /api/users                 & Oui & Non & Non \\
GET  & /api/pfe/sujets            & Oui & Oui & Oui \\
POST & /api/pfe/sujets            & Oui & Oui & Non \\
GET  & /api/requests/admin/inbox  & Oui & Non & Non \\
POST & /api/requests/reclamations & Non & Non & Oui \\
\bottomrule
\end{tabular}
\end{table}

\section{Implementation Backend}

\subsection{Service d'Authentification}

Le service encapsule toute la logique liee a la verification des
identifiants, au hachage des mots de passe et a la generation des
tokens JWT.

Les principales fonctions implementees sont :
\begin{itemize}
  \item \texttt{loginUser(email, password)} : verifie les identifiants
        et retourne un token JWT valide 24 heures.
  \item \texttt{verifyToken(token)} : decode et valide la signature
        d'un token recu.
  \item \texttt{createUser(email, password, role)} : hache le mot de
        passe avec bcrypt (cout 12) et cree le compte.
\end{itemize}

\subsection{Middleware d'Authentification}

Le middleware \texttt{requireAuth} verifie la presence et la validite
du token JWT dans l'en-tete \texttt{Authorization} de chaque requete
entrante. En cas d'echec, il retourne un code 401.

\subsection{Middleware d'Autorisation par Role}

Le middleware \texttt{requireRole} est une fonction d'ordre superieur
qui genere dynamiquement un middleware verifiant que le role de
l'utilisateur fait partie de la liste des roles autorises pour la
route concernee. En cas d'echec, il retourne un code 403.

\subsection{Controleur d'Authentification}

Le controleur expose trois routes principales :
\begin{itemize}
  \item \texttt{POST /api/auth/login} : connexion utilisateur.
  \item \texttt{GET /api/auth/me} : recuperation du profil de
        l'utilisateur connecte (route protegee).
  \item \texttt{POST /api/auth/logout} : deconnexion (invalidation
        cote client).
\end{itemize}

\section{Implementation Frontend}

\subsection{Contexte d'Authentification React}

Le \texttt{AuthContext} centralise l'etat d'authentification et le
rend disponible a tous les composants de l'application via le hook
\texttt{useAuth()}. Il expose :
\begin{itemize}
  \item L'objet \texttt{user} (ou \texttt{null}).
  \item La fonction \texttt{login(email, password)}.
  \item La fonction \texttt{logout()}.
  \item L'etat \texttt{loading} pour la verification initiale.
\end{itemize}

\subsection{Intercepteur Axios}

L'intercepteur injecte automatiquement le token JWT dans l'en-tete de
chaque requete sortante et redirige vers la page de connexion en cas
de reponse 401.

\subsection{Routes Protegees}

Le composant \texttt{ProtectedRoute} empeche l'acces aux pages privees
si l'utilisateur n'est pas authentifie ou ne possede pas le role
requis. La redirection se fait automatiquement vers la page de
connexion ou la page d'erreur 403.

\section{Securite Applicative}

\begin{table}[H]
\centering
\caption{Mesures de securite implementees}
\begin{tabular}{lp{8cm}}
\toprule
\textbf{Menace} & \textbf{Contre-mesure} \\
\midrule
Injection SQL    & Prisma ORM avec requetes parametrees \\
Force brute      & Rate-limiting 20 requetes / 15 min par IP \\
XSS              & React echappe le HTML par defaut, en-tetes CSP \\
CSRF             & JWT en en-tete (non envoye automatiquement) \\
JWT falsifie     & Signature HMAC-SHA256, cle de 256 bits minimum \\
Fuite de donnees & select Prisma explicite, jamais retourner le hash \\
Elevation role   & Middleware requireRole sur chaque route sensible \\
\bottomrule
\end{tabular}
\end{table}

\section{Tests et Validation}

\subsection{Tests Fonctionnels -- Authentification}

\begin{table}[H]
\centering
\caption{Resultats des tests fonctionnels d'authentification}
\begin{tabular}{lp{6cm}cc}
\toprule
\textbf{ID} & \textbf{Scenario} & \textbf{Code attendu} & \textbf{Resultat} \\
\midrule
TF-01 & Connexion identifiants valides         & 200 & PASS \\
TF-02 & Connexion mot de passe incorrect       & 401 & PASS \\
TF-03 & Connexion email inexistant             & 401 & PASS \\
TF-04 & Acces route protegee sans token        & 401 & PASS \\
TF-05 & Acces route protegee avec token expire & 401 & PASS \\
TF-06 & Acces route admin avec role enseignant & 403 & PASS \\
TF-07 & Recuperation profil /auth/me           & 200 & PASS \\
TF-08 & Deconnexion utilisateur                & 200 & PASS \\
\bottomrule
\end{tabular}
\end{table}

\subsection{Tests Fonctionnels -- RBAC}

\begin{table}[H]
\centering
\caption{Resultats des tests RBAC}
\begin{tabular}{lp{6cm}cc}
\toprule
\textbf{ID} & \textbf{Scenario} & \textbf{Code attendu} & \textbf{Resultat} \\
\midrule
TC-01 & Admin cree un compte enseignant         & 201 & PASS \\
TC-02 & Enseignant tente de creer un compte     & 403 & PASS \\
TC-03 & Etudiant soumet une reclamation         & 201 & PASS \\
TC-04 & Etudiant accede a l'inbox admin         & 403 & PASS \\
TC-05 & Enseignant modifie sujet PFE (propose)  & 200 & PASS \\
TC-06 & Enseignant modifie sujet PFE valide     & 403 & PASS \\
TC-07 & Admin modifie tout sujet PFE            & 200 & PASS \\
\bottomrule
\end{tabular}
\end{table}

\subsection{Tests de Securite}

\begin{table}[H]
\centering
\caption{Resultats des tests de securite}
\begin{tabular}{lp{6cm}c}
\toprule
\textbf{Test} & \textbf{Description} & \textbf{Resultat} \\
\midrule
Brute-force /login &
  50 requetes en 10 secondes &
  429 apres 20 req \\
JWT falsifie &
  Modification d'un bit dans la signature &
  401 PASS \\
Elevation privilege &
  Token student sur route admin &
  403 PASS \\
Injection SQL &
  Payload OR 1=1 dans le champ email &
  Neutralise PASS \\
\bottomrule
\end{tabular}
\end{table}

\subsection{Couverture de Tests}

\begin{table}[H]
\centering
\caption{Couverture de tests par module}
\begin{tabular}{lc}
\toprule
\textbf{Module} & \textbf{Couverture (\%)} \\
\midrule
Auth         & 92 \\
Users        & 88 \\
PFE          & 85 \\
Requests     & 90 \\
Annonces     & 78 \\
\midrule
\textbf{Moyenne} & \textbf{86.6} \\
\bottomrule
\end{tabular}
\end{table}

\section{Conclusion}

Ce chapitre a presente l'implementation complete du module
d'authentification et de gestion des roles, couvrant le backend
(service JWT, middlewares RBAC, controleur) et le frontend (contexte
React, intercepteur Axios, routes protegees). La campagne de tests
fonctionnels et de securite confirme la robustesse du systeme face aux
principales menaces applicatives.

% ════════════════════════════════════════════════════════════
%  CONCLUSION GENERALE
% ════════════════════════════════════════════════════════════
\chapter*{Conclusion Generale et Perspectives}
\addcontentsline{toc}{chapter}{Conclusion Generale}

\section*{Bilan}

Ce memoire a presente la conception et le developpement du module
\textbf{Authentification et Gestion des Roles} du Systeme de Gestion
Universitaire Integre (SGUI). Le travail realise a permis d'atteindre
l'ensemble des objectifs fixes :

\begin{itemize}
  \item Un mecanisme d'authentification JWT conforme au standard
        RFC 7519, garantissant des sessions sans etat et une
        scalabilite horizontale.
  \item Un modele RBAC a trois roles (Administrateur, Enseignant,
        Etudiant) avec une matrice d'autorisation fine par route API.
  \item Une interface utilisateur reactive et moderne, construite avec
        React.js et Tailwind CSS, offrant des tableaux de bord adaptes
        a chaque profil.
  \item Une securite applicative validee contre les principales
        vulnerabilites web (injection SQL, XSS, CSRF, force brute,
        elevation de privileges).
  \item Une couverture de tests superieure a 85\% sur l'ensemble des
        modules critiques.
\end{itemize}

\section*{Difficultes Rencontrees}

Le principal defi technique a ete la gestion de la synchronisation des
roles entre le token JWT (contenant le role au moment de la connexion)
et les modifications de role effectuees par l'administrateur apres
l'emission du token. Ce probleme a ete resolu par l'implementation d'un
mecanisme de verification en base de donnees sur les routes critiques.

\section*{Perspectives}

Plusieurs ameliorations peuvent etre envisagees pour les versions
futures :

\begin{enumerate}
  \item \textbf{Authentification multi-facteurs (MFA)} : ajout d'un
        second facteur (TOTP, SMS) pour les comptes administrateurs.
  \item \textbf{Refresh Tokens} : mecanisme de renouvellement silencieux
        des tokens JWT pour ameliorer l'experience utilisateur.
  \item \textbf{Single Sign-On (SSO)} : integration avec le systeme
        d'identite de l'universite via OpenID Connect ou SAML 2.0.
  \item \textbf{Audit Log} : journalisation immuable de toutes les
        actions sensibles (connexion, modification de role, validation
        PFE).
  \item \textbf{Internationalisation} : prise en charge complete de
        l'arabe pour les interfaces utilisateur.
\end{enumerate}

\vspace{1cm}

En conclusion, ce projet a demontre la viabilite d'une architecture
moderne basee sur JWT et RBAC pour la gestion securisee des acces dans
un contexte universitaire algerien. Le systeme est deployable,
maintenable et extensible, constituant une base solide pour les futurs
modules du SGUI.

% ────────────────────────────────────────────────────────────
%  REFERENCES
% ────────────────────────────────────────────────────────────
\chapter*{References Bibliographiques}
\addcontentsline{toc}{chapter}{References Bibliographiques}

\begin{enumerate}
  \item M. Jones, J. Bradley, N. Sakimura, ``JSON Web Token (JWT)'',
        RFC 7519, IETF, 2015.

  \item M. Jones, D. Hardt, ``The OAuth 2.0 Authorization Framework:
        Bearer Token Usage'', RFC 6750, IETF, 2012.

  \item OWASP Foundation, ``OWASP Top Ten 2021'', OWASP, 2021.

  \item M. Casciaro, L. Mammino, \textit{Node.js Design Patterns},
        3rd edition, Packt Publishing, 2020.

  \item Prisma Team, \textit{Prisma Documentation}, 2024.\\
        \url{https://www.prisma.io/docs}

  \item N. Provos, D. Mazieres, ``A Future-Adaptable Password Scheme'',
        USENIX Annual Technical Conference, 1999.

  \item D. Ferraiolo, R. Kuhn, R. Chandramouli, \textit{Role-Based
        Access Control}, Artech House, 2003.

  \item R. Wieruch, \textit{The Road to React}, 2023.

  \item Express.js Documentation, \url{https://expressjs.com}

  \item React Documentation, \url{https://react.dev}
\end{enumerate}

% ============================================================
\end{document}
% ============================================================
